\documentclass[border=10pt]{standalone}
\usepackage[utf8]{inputenc}
\usepackage[T1]{fontenc}
\usepackage{tikz}
\usetikzlibrary{positioning, arrows.meta, decorations.pathmorphing, calc, backgrounds}

\begin{document}
\begin{tikzpicture}[
    arrow/.style={
        -{Stealth[length=6pt, width=5pt]}, line width=1.2pt, color=#1
    },
    label box/.style={
        font=\sffamily\small, text=black, align=left, text width=7cm
    },
    bold label/.style={
        font=\sffamily\small\bfseries, text=#1
    }
]

% ── Colores ──
\definecolor{mainblue}{HTML}{0D47A1}
\definecolor{accentteal}{HTML}{00897B}
\definecolor{rawgray}{HTML}{616161}
\definecolor{fftblue}{HTML}{1565C0}
\definecolor{cwtgreen}{HTML}{2E7D32}
\definecolor{wstpurp}{HTML}{5C6BC0}
\definecolor{momorange}{HTML}{F57C00}
\definecolor{bglight}{HTML}{F5F7FA}

% ════════════════════════════════════════
% TÍTULO
% ════════════════════════════════════════
\node[font=\sffamily\Large\bfseries, text=mainblue, align=center]
    at (6, 1.5) {La Búsqueda de la Mejor Representación};
\node[font=\sffamily\small, text=black!60, align=center]
    at (6, 0.8) {¿Cómo debe observar la señal el algoritmo? Evaluamos 5 transformaciones fundamentales.};

% ════════════════════════════════════════
% SEÑAL DE ENTRADA (onda sinusoidal)
% ════════════════════════════════════════
\begin{scope}[shift={(-1.5, -2.5)}]
    % Caja de fondo
    \fill[mainblue!8, rounded corners=6pt] (-1.8, -1.2) rectangle (1.8, 1.2);
    \draw[mainblue!40, rounded corners=6pt, line width=0.8pt] (-1.8, -1.2) rectangle (1.8, 1.2);

    % Onda sísmica
    \draw[mainblue, line width=1.5pt, smooth, samples=80, domain=-1.4:1.4]
        plot (\x, {0.7*sin(900*\x)*exp(-abs(\x)*1.2)});
    \draw[accentteal, line width=1pt, smooth, samples=80, domain=-1.4:1.4]
        plot (\x, {0.4*sin(1200*\x+60)*exp(-abs(\x)*1.5) - 0.1});

    \node[font=\sffamily\tiny\bfseries, text=mainblue] at (0, -0.9) {Señal S1 + S2};
\end{scope}

% ════════════════════════════════════════
% PRISMA (triángulo que "descompone")
% ════════════════════════════════════════
\begin{scope}[shift={(2.2, -2.5)}]
    % Prisma triangular con gradiente simulado
    \fill[mainblue!15] (0, 0.9) -- (-0.7, -0.7) -- (0.7, -0.7) -- cycle;
    \fill[mainblue!25] (0, 0.9) -- (0.7, -0.7) -- (1.0, -0.4) -- (0.3, 1.2) -- cycle;
    \draw[mainblue!60, line width=0.8pt] (0, 0.9) -- (-0.7, -0.7) -- (0.7, -0.7) -- cycle;
    \draw[mainblue!60, line width=0.8pt] (0, 0.9) -- (0.3, 1.2) -- (1.0, -0.4) -- (0.7, -0.7);
\end{scope}

% Flecha entrada → prisma
\draw[arrow=mainblue] (0.1, -2.5) -- (1.3, -2.5);

% ════════════════════════════════════════
% 5 FLECHAS divergentes desde el prisma
% ════════════════════════════════════════
\draw[arrow=rawgray]  (3.1, -1.8) -- (4.8, -0.5);
\draw[arrow=fftblue]  (3.2, -2.1) -- (4.8, -1.7);
\draw[arrow=cwtgreen]  (3.2, -2.5) -- (4.8, -2.9);
\draw[arrow=wstpurp]  (3.2, -2.9) -- (4.8, -4.1);
\draw[arrow=momorange] (3.1, -3.2) -- (4.8, -5.3);

% ════════════════════════════════════════
% REPRESENTACIÓN 1: Raw (barras de voltaje)
% ════════════════════════════════════════
\begin{scope}[shift={(5.5, -0.5)}]
    \fill[rawgray!8, rounded corners=3pt] (-0.5, -0.45) rectangle (0.5, 0.45);
    \draw[rawgray!40, rounded corners=3pt, line width=0.5pt] (-0.5, -0.45) rectangle (0.5, 0.45);
    % Barras irregulares (voltaje crudo)
    \fill[rawgray!70] (-0.35, -0.3) rectangle (-0.25, -0.1);
    \fill[rawgray]     (-0.2, -0.3) rectangle (-0.1, 0.25);
    \fill[rawgray!80] (-0.05, -0.3) rectangle (0.05, 0.1);
    \fill[rawgray!50] (0.1, -0.3)  rectangle (0.2, -0.05);
    \fill[rawgray!60] (0.25, -0.3) rectangle (0.35, 0.0);
\end{scope}
\node[label box, anchor=west] at (6.3, -0.5) {
    \textcolor{rawgray}{\textbf{Voltajes Crudos (Raw):}} El voltaje
    temporal puro sin alteración.
};

% ════════════════════════════════════════
% REPRESENTACIÓN 2: FFT (onda sinusoidal)
% ════════════════════════════════════════
\begin{scope}[shift={(5.5, -1.7)}]
    \fill[fftblue!8, rounded corners=3pt] (-0.5, -0.45) rectangle (0.5, 0.45);
    \draw[fftblue!40, rounded corners=3pt, line width=0.5pt] (-0.5, -0.45) rectangle (0.5, 0.45);
    \draw[fftblue, line width=1pt, smooth, samples=40, domain=-0.35:0.35]
        plot (\x, {0.3*sin(1400*\x)});
\end{scope}
\node[label box, anchor=west] at (6.3, -1.7) {
    \textcolor{fftblue}{\textbf{Transformada de Fourier (FFT):}} Espectros
    de magnitud en el dominio de frecuencia.
};

% ════════════════════════════════════════
% REPRESENTACIÓN 3: CWT (escalograma — grid de colores)
% ════════════════════════════════════════
\begin{scope}[shift={(5.5, -2.9)}]
    \fill[cwtgreen!8, rounded corners=3pt] (-0.5, -0.45) rectangle (0.5, 0.45);
    \draw[cwtgreen!40, rounded corners=3pt, line width=0.5pt] (-0.5, -0.45) rectangle (0.5, 0.45);
    % Mini escalograma (grid de colores)
    \foreach \x/\c in {-0.38/red!60, -0.24/orange!70, -0.10/yellow!80, 0.04/green!60, 0.18/cyan!60, 0.32/blue!50} {
        \foreach \y/\o in {-0.28/40, -0.14/70, 0.0/100, 0.14/60, 0.28/30} {
            \fill[\c!\o] (\x, \y) rectangle ++(0.12, 0.12);
        }
    }
\end{scope}
\node[label box, anchor=west] at (6.3, -2.9) {
    \textcolor{cwtgreen}{\textbf{Wavelet Continua (CWT):}} Escalogramas
    de tiempo y frecuencia.
};

% ════════════════════════════════════════
% REPRESENTACIÓN 4: WST (coeficientes — bloques escalonados)
% ════════════════════════════════════════
\begin{scope}[shift={(5.5, -4.1)}]
    \fill[wstpurp!8, rounded corners=3pt] (-0.5, -0.45) rectangle (0.5, 0.45);
    \draw[wstpurp!40, rounded corners=3pt, line width=0.5pt] (-0.5, -0.45) rectangle (0.5, 0.45);
    % Bloques escalonados (scattering orders)
    \fill[wstpurp!30] (-0.35, -0.25) rectangle (-0.05, -0.05);
    \fill[wstpurp!50] (-0.35, 0.0)  rectangle (0.15, 0.2);
    \fill[wstpurp!70] (-0.35, 0.22) rectangle (0.35, 0.35);
    % Líneas de subdivisión
    \draw[wstpurp!40, line width=0.3pt] (-0.2, -0.25) -- (-0.2, -0.05);
    \draw[wstpurp!40, line width=0.3pt] (-0.1, 0.0) -- (-0.1, 0.2);
    \draw[wstpurp!40, line width=0.3pt] (0.0, 0.0) -- (0.0, 0.2);
    \draw[wstpurp!40, line width=0.3pt] (-0.15, 0.22) -- (-0.15, 0.35);
    \draw[wstpurp!40, line width=0.3pt] (0.05, 0.22) -- (0.05, 0.35);
    \draw[wstpurp!40, line width=0.3pt] (0.2, 0.22) -- (0.2, 0.35);
\end{scope}
\node[label box, anchor=west] at (6.3, -4.1) {
    \textcolor{wstpurp}{\textbf{Wavelet Scattering (WST):}} Coeficientes
    compactos estables a deformaciones.
};

% ════════════════════════════════════════
% REPRESENTACIÓN 5: MOMENT (red neuronal — nodos)
% ════════════════════════════════════════
\begin{scope}[shift={(5.5, -5.3)}]
    \fill[momorange!8, rounded corners=3pt] (-0.5, -0.45) rectangle (0.5, 0.45);
    \draw[momorange!40, rounded corners=3pt, line width=0.5pt] (-0.5, -0.45) rectangle (0.5, 0.45);
    % Red neuronal simplificada (3 capas)
    \foreach \y in {-0.2, 0, 0.2} {
        \fill[momorange!60] (-0.3, \y) circle (0.06);
    }
    \foreach \y in {-0.25, -0.05, 0.15, 0.35} {
        \fill[momorange!80] (0, \y-0.05) circle (0.06);
    }
    \foreach \y in {-0.15, 0.05, 0.25} {
        \fill[momorange] (0.3, \y) circle (0.06);
    }
    % Conexiones
    \foreach \a in {-0.2, 0, 0.2} {
        \foreach \b in {-0.25, -0.05, 0.15, 0.35} {
            \draw[momorange!30, line width=0.2pt] (-0.3, \a) -- (0, \b-0.05);
        }
    }
    \foreach \a in {-0.25, -0.05, 0.15, 0.35} {
        \foreach \b in {-0.15, 0.05, 0.25} {
            \draw[momorange!30, line width=0.2pt] (0, \a-0.05) -- (0.3, \b);
        }
    }
\end{scope}
\node[label box, anchor=west] at (6.3, -5.3) {
    \textcolor{momorange}{\textbf{Modelo Fundacional (MOMENT):}} Embeddings
    genéricos de un modelo pre-entrenado.
};

% ════════════════════════════════════════
% FONDO GENERAL
% ════════════════════════════════════════
\begin{scope}[on background layer]
    \fill[bglight, rounded corners=10pt] (-3.8, 2.2) rectangle (14, -6.2);
\end{scope}

\end{tikzpicture}
\end{document}
